%%%%%%%%%%%%%%%%%%%%%%%%%%%%%%%%%%%%%%%%%%%%%%%%%%%%%%%%%%%%%%%%
%  chapter3.tex  --  Avaliacao Experimental
%  Jogo ISOLA em Prolog com Algoritmo Alfa-Beta
%  Inteligencia Artificial, ISEL 2025/2026
%%%%%%%%%%%%%%%%%%%%%%%%%%%%%%%%%%%%%%%%%%%%%%%%%%%%%%%%%%%%%%%%

\chapter{Avalia\c{c}\~{a}o Experimental}
\label{chap:eval}

Este cap\'{i}tulo apresenta os resultados dos testes realizados ao sistema:
testes unit\'{a}rios de correc\c{c}\~{a}o funcional e benchmarks de desempenho
em diferentes configura\c{c}\~{o}es de tabuleiro e profundidade de pesquisa.

% ===================================================================
\section{Testes Unit\'{a}rios}
\label{sec:unit_tests}

O projeto inclui uma suite de 33 testes unit\'{a}rios (ficheiro
\texttt{tests/tests.pl}), organizados em quatro grupos.

\subsection{Execu\c{c}\~{a}o e Resultados}

Os testes foram executados com o comando:

\begin{verbatim}
swipl -g "run_tests, halt" tests/tests.pl
\end{verbatim}

O resultado obtido foi:

\begin{verbatim}
% Start unit: board
% [1/33] board:initial_board_6x6_size ............... passed (0.015 sec)
% [2/33] board:initial_board_6x6_row_length ......... passed (0.001 sec)
% [3/33] board:initial_board_all_free ............... passed (0.000 sec)
% [4/33] board:all_cells_3x3_count .................. passed (0.000 sec)
% [5/33] board:get_cell_after_set ................... passed (0.000 sec)
% [6/33] board:set_cell_nondestructive .............. passed (0.000 sec)
% [7/33] board:get_cell_oob_row ..................... passed (0.000 sec)
% [8/33] board:get_cell_oob_col ..................... passed (0.000 sec)
% [9/33] board:empty_cell_free ...................... passed (0.000 sec)
% [10/33] board:empty_cell_pawn ..................... passed (0.000 sec)
% [11/33] board:empty_cell_removed .................. passed (0.000 sec)
% End unit board: passed (0.007 sec CPU)

% Start unit: rules
% [12/33] rules:valid_moves_centre_3x3_count ........ passed (0.000 sec)
% [13/33] rules:valid_moves_corner_3x3_count ........ passed (0.000 sec)
% [14/33] rules:valid_moves_blocked_by_remove ....... passed (0.000 sec)
% [15/33] rules:valid_moves_no_overlap_opponent ..... passed (0.000 sec)
% [16/33] rules:valid_removes_excludes_pawns ........ passed (0.000 sec)
% [17/33] rules:valid_removes_excl_removed_cell ..... passed (0.000 sec)
% [18/33] rules:apply_move_source_cleared ........... passed (0.000 sec)
% [19/33] rules:apply_remove_marks_cell ............. passed (0.000 sec)
% [20/33] rules:game_over_no_moves .................. passed (0.000 sec)
% [21/33] rules:game_over_has_moves ................. passed (0.000 sec)
% [22/33] rules:player_pos_p1 ....................... passed (0.000 sec)
% [23/33] rules:player_pos_p2 ....................... passed (0.000 sec)
% End unit rules: passed (0.001 sec CPU)

% Start unit: ai
% [24/33] ai:evaluate_mobility_advantage ............ passed (0.000 sec)
% [25/33] ai:evaluate_opponent_advantage ............ passed (0.000 sec)
% [26/33] ai:evaluate_symmetric ..................... passed (0.000 sec)
% [27/33] ai:terminal_value_p1_loses ................ passed (0.000 sec)
% [28/33] ai:terminal_value_p2_wins ................. passed (0.000 sec)
% [29/33] ai:best_move_trivial_win .................. passed (0.001 sec)
% [30/33] ai:alphabeta_terminal ..................... passed (0.000 sec)
% [31/33] ai:minimax_depth_zero ..................... passed (0.000 sec)
% End unit ai: passed (0.002 sec CPU)

% Start unit: integration
% [32/33] integration:scripted_turn_1 ............... passed (0.000 sec)
% [33/33] integration:scripted_game_over_detection .. passed (0.000 sec)
% End unit integration: passed (0.000 sec CPU)

% All 33 tests passed in 0.029 seconds (0.014 cpu)
\end{verbatim}

Todos os 33 testes passaram em 0.029~segundos. A execu\c{c}\~{a}o r\'{a}pida
deve-se ao facto de a maioria dos testes operar sobre tabuleiros 3{\texttimes}3
ou posi\c{c}\~{o}es artificiais constru\'{i}das especificamente para o teste.

\subsection{Cobertura dos Testes}

\begin{table}[htbp]
  \caption{Distribui\c{c}\~{a}o e cobertura dos testes unit\'{a}rios}
  \label{tab:tests}
  \centering
  \begin{tabular}{llp{7cm}}
    \toprule
    \textbf{M\'{o}dulo} & \textbf{N\'{u}m. testes} & \textbf{Aspetos verificados} \\
    \midrule
    \texttt{board}       & 11 &
      Cria\c{c}\~{a}o do tabuleiro; acesso por coordenadas;
      atualiza\c{c}\~{a}o n\~{a}o destrutiva; acesso fora dos limites
      (deve falhar sem erro); \texttt{empty\_cell} para
      c\'{e}lulas livres, ocupadas e removidas. \\
    \addlinespace
    \texttt{rules}       & 12 &
      Movimentos v\'{a}lidos do centro e do canto;
      bloqueio por casa removida; aus\^{e}ncia de sobreposic\~{a}o
      com o advers\'{a}rio; lista de remo\c{c}\~{o}es (exclu\'{i} pe\~{o}es
      e casas j\'{a} removidas); aplica\c{c}\~{a}o de movimento e
      remo\c{c}\~{a}o; detec\c{c}\~{a}o de fim de jogo. \\
    \addlinespace
    \texttt{ai}          & 8 &
      Fun\c{c}\~{a}o de avalia\c{c}\~{a}o (vantagem P1, vantagem P2,
      posi\c{c}\~{a}o sim\'{e}trica); valores terminais; melhor
      movimento em situa\c{c}\~{a}o de vit\'{o}ria trivial;
      alfa-beta em posi\c{c}\~{a}o terminal; minimax a profundidade zero. \\
    \addlinespace
    \texttt{integration} & 2 &
      Turno completo (movimento + remo\c{c}\~{a}o num jogo real);
      detec\c{c}\~{a}o de fim de jogo no contexto de um jogo jogado. \\
    \bottomrule
  \end{tabular}
\end{table}

\subsection{Exemplo de Teste Unit\'{a}rio}

Para ilustrar a estrutura dos testes, apresenta-se o teste que verifica
o n\'{u}mero de movimentos v\'{a}lidos a partir do centro de um tabuleiro 3{\texttimes}3:

\begin{lstlisting}[language=Prolog,
    caption={Teste: movimentos do centro de um tabuleiro 3{\texttimes}3},
    label={lst:test_example},
    numbers=left,
    basicstyle=\small\ttfamily,
    frame=single]
:- use_module(library(plunit)).
:- begin_tests(rules).

test(valid_moves_centre_3x3_count) :-
    initial_board(3, B0),
    set_cell(B0, 0, 0, x, B1),
    set_cell(B1, 2, 2, o, Board),
    valid_moves(Board, 1, 1, Moves), % piao no centro (1,1)
    length(Moves, N),
    assertion(N =:= 7). % 8 direcoes - 1 (canto ocupado)

:- end_tests(rules).
\end{lstlisting}

% ===================================================================
\section{Benchmarks de Desempenho}
\label{sec:benchmarks}

Para avaliar o desempenho do algoritmo alfa-beta, foram medidos os
tempos de execu\c{c}\~{a}o do primeiro movimento da IA em diferentes
combina\c{c}\~{o}es de tamanho de tabuleiro e profundidade de pesquisa.
Todos os testes partiram do estado inicial do jogo (pe\~{o}es nos cantos
opostos, tabuleiro completamente livre).

\subsection{Metodologia}

Os tempos foram medidos com o predicado \texttt{get\_time/1} do
SWI-Prolog (tempo de rel\'{o}gio de parede), no seguinte ambiente:

\begin{itemize}
  \item Sistema operativo: macOS Darwin 25.3.0 (Apple Silicon)
  \item SWI-Prolog (vers\~{a}o padr\~{a}o, sem otimiza\c{c}\~{o}es adicionais)
  \item Medida: tempo para calcular o primeiro movimento da IA
        a partir do estado inicial
\end{itemize}

\subsection{Resultados}

\begin{table}[htbp]
  \caption{Tempo de c\'{a}lculo do primeiro movimento da IA}
  \label{tab:benchmarks}
  \centering
  \begin{tabular}{ccrll}
    \toprule
    \textbf{Tabuleiro} &
    \textbf{Profundidade} &
    \textbf{Tempo} &
    \textbf{Movimento} &
    \textbf{Remo\c{c}\~{a}o} \\
    \midrule
    3{\texttimes}3 & 2 & 4~ms       & (0,0)$\to$(1,1) & (1,2) \\
    3{\texttimes}3 & 4 & 136~ms     & (0,0)$\to$(1,1) & (1,2) \\
    6{\texttimes}6 & 2 & 87~ms      & (0,0)$\to$(1,1) & (4,4) \\
    6{\texttimes}6 & 4 & 167.619~ms $\approx$ 2 min 47 s & (0,0)$\to$(1,1) & (4,4) \\
    \bottomrule
  \end{tabular}
\end{table}

\subsection{An\'{a}lise dos Resultados}

\subsubsection{Fator de Ramifica\c{c}\~{a}o Real}

No estado inicial de um tabuleiro 6{\texttimes}6, o Jogador~1 pode:
\begin{itemize}
  \item \textbf{Mover} para 3 casas (canto com 3 vizinhos livres).
  \item \textbf{Remover} qualquer uma de 34 casas livres
        (36 casas $-$ 2 pe\~{o}es).
\end{itemize}

Total de pares (movimento, remo\c{c}\~{a}o): $3 \times 34 = 102$.

Este \'{e} o fator de ramifica\c{c}\~{a}o no n\'{i}vel 1. Nas jogadas seguintes,
quando os pe\~{o}es se movem para o centro, cada jogador pode ter at\'{e}~8
destinos poss\'{i}veis, o que eleva o fator de ramifica\c{c}\~{a}o.

\subsubsection{Impacto da Profundidade}

O salto de profundidade~2 para profundidade~4 no tabuleiro 6{\texttimes}6
representa um aumento de $87~\text{ms}$ para $167.619~\text{ms}$ ---
um fator de aproximadamente $\mathbf{1927\times}$. Este crescimento
super-exponencial confirma o car\'{a}ter combinat\'{o}rio do problema.

Embora o alfa-beta reduza significativamente o n\'{u}mero de n\'{o}s
explorados em rela\c{c}\~{a}o ao minimax puro, o fator de ramifica\c{c}\~{a}o
elevado ($\approx 102$) no estado inicial de um tabuleiro 6{\texttimes}6
torna a profundidade~4 impraticável para uso interativo.

\subsubsection{Consist\^{e}ncia dos Movimentos Escolhidos}

Nota-se que em todas as configura\c{c}\~{o}es, a IA escolheu o mesmo
movimento: mover o pe\~{a}o do canto (0,0) para o centro (1,1).
Isto \'{e} consistente com a fun\c{c}\~{a}o de avalia\c{c}\~{a}o que favorece
a centralidade --- o pe\~{a}o no centro (1,1) de um tabuleiro 3{\texttimes}3
tem acesso a at\'{e}~7 casas, enquanto no canto (0,0) tem apenas~3.

A diferen\c{c}a na escolha da remo\c{c}\~{a}o entre profundidade~2
((1,2) em 3{\texttimes}3) e profundidade~4 ((4,4) em 6{\texttimes}6)
reflete a vis\~{a}o de futuro mais profunda da pesquisa a maior
profundidade.

\subsubsection{Avalia\c{c}\~{a}o da Fun\c{c}\~{a}o Heur\'{i}stica}

Os testes confirm\'{a}ram os valores esperados da fun\c{c}\~{a}o de avalia\c{c}\~{a}o:

\begin{table}[htbp]
  \caption{Valores da fun\c{c}\~{a}o de avalia\c{c}\~{a}o em posi\c{c}\~{o}es de teste}
  \label{tab:eval_values}
  \centering
  \begin{tabular}{llrrr}
    \toprule
    \textbf{Posi\c{c}\~{a}o} &
    \textbf{Configura\c{c}\~{a}o} &
    \textbf{Mob. P1} &
    \textbf{Mob. P2} &
    \textbf{Valor} \\
    \midrule
    Sim\'{e}trica (cantos opostos)  & 3{\texttimes}3, (0,0) vs.\ (2,2) & 3 & 3 & 0 \\
    Assim\'{e}trica (canto vs.\ centro) & 3{\texttimes}3, (0,0) vs.\ (1,1) & 2 & 7 & $-52$ \\
    \bottomrule
  \end{tabular}
\end{table}

O valor~$0$ na posi\c{c}\~{a}o sim\'{e}trica confirma que a fun\c{c}\~{a}o n\~{a}o \'{e}
tendenciosa para nenhum jogador. O valor~$-52$ na posi\c{c}\~{a}o assim\'{e}trica
reflete corretamente a vantagem significativa de P2 (maior mobilidade e
melhor centralidade).

% ===================================================================
\section{Limita\c{c}\~{o}es e Trabalho Futuro}
\label{sec:limitations}

\subsection{Limita\c{c}\~{o}es Identificadas}

\begin{enumerate}
  \item \textbf{Desempenho em tabuleiros grandes a profundidade alta:}
        A combina\c{c}\~{a}o 6{\texttimes}6 + profundidade~4 \'{e} proibitiva
        ($\approx 167$~s por movimento). O fator de ramifica\c{c}\~{a}o
        elevado, resultante da grande quantidade de casas remov\'{i}veis,
        \'{e} a causa principal.

  \item \textbf{Formato de entrada:} os jogadores humanos t\^{e}m de
        introduzir posi\c{c}\~{o}es no formato \texttt{Linha/Coluna.}
        (com ponto final), pois o predicado \texttt{read/1} do
        SWI-Prolog espera um termo Prolog completo terminado por ponto.
        Este formato pode causar confus\~{a}o.

  \item \textbf{Aus\^{e}ncia de tabela de transposi\c{c}\~{a}o:} posi\c{c}\~{o}es
        iguais atingidas por caminhos diferentes s\~{a}o reavaliadas.
        Uma tabela de transposi\c{c}\~{a}o reduziria o trabalho redundante.

  \item \textbf{Sem aprofundamento iterativo:} a profundidade \'{e}
        fixa. O aprofundamento iterativo (aumentar gradualmente a
        profundidade at\'{e} ao limite de tempo) melhoraria a qualidade
        do jogo em tempo fixo.
\end{enumerate}

\subsection{Recomenda\c{c}\~{o}es de Configura\c{c}\~{a}o}

Com base nos resultados experimentais:

\begin{itemize}
  \item Para tabuleiros 3{\texttimes}3: profundidade~4 \'{e} adequada
        (136~ms por movimento).
  \item Para tabuleiros 6{\texttimes}6: profundidade~2 \'{e} o
        compromisso recomendado (87~ms por movimento),
        oferecendo uma resposta quase instant\^{a}nea.
  \item Profundidade~6 n\~{a}o \'{e} recomendada para tabuleiros
        maiores que 3{\texttimes}3 sem otimiza\c{c}\~{o}es adicionais.
\end{itemize}

% ===================================================================
\section{Conclus\~{o}es}
\label{sec:conclusions}

Este trabalho demonstrou que \'{e} poss\'{i}vel implementar um jogo de
estrat\'{e}gia completo com IA em Prolog de forma elegante e eficaz.

\textbf{Objetivos cumpridos:}
\begin{itemize}
  \item Implementa\c{c}\~{a}o completa do jogo ISOLA com todos os modos
        (humano-humano, humano-IA, IA-IA) e tamanhos de tabuleiro
        configur\'{a}veis.
  \item Algoritmo alfa-beta funcional com ordena\c{c}\~{a}o de movimentos
        e fun\c{c}\~{a}o de avalia\c{c}\~{a}o heur\'{i}stica.
  \item Suite de 33 testes unit\'{a}rios, todos bem sucedidos.
  \item Compreens\~{a}o pr\'{a}tica do \textit{backtracking}, da programa\c{c}\~{a}o
        l\'{o}gica e dos algoritmos de pesquisa em jogos.
\end{itemize}

\textbf{Principais aprendizagens:}
\begin{itemize}
  \item O Prolog permite exprimir regras de jogo de forma quase
        direta a partir da descri\c{c}\~{a}o em linguagem natural,
        aproveitando o \textit{backtracking} para gerar automaticamente
        todos os movimentos v\'{a}lidos.
  \item O alfa-beta \'{e} essencial para tornar o \textit{minimax}
        pr\'{a}tico: sem poda, profundidade~4 num tabuleiro 6{\texttimes}6
        seria completamente invi\'{a}vel.
  \item A ordena\c{c}\~{a}o de movimentos \'{e} t\~{a}o importante quanto o
        pr\'{o}prio alfa-beta: maus movimentos avaliados primeiro
        reduzem drasticamente a efic\'{a}cia dos cortes.
\end{itemize}
